% Template for the PDF version of a Discrete Mathematics and Algorithms exam.
% Machine translated from the Danish original, ../danish/template.tex
% Copy this file (e.g. to reexam-april-2026.tex) and adjust the settings below.
% Remove the % in front of \answerstrue for the version with answers and explanations.
% Every \placeholder{...} describes information to fill in for a specific exam.
\documentclass[12pt]{article}
\newif\ifanswers

% ===== Settings for a specific exam =====
%\answerstrue % with answers and explanations
\newcommand{\placeholder}[1]{\textit{[#1]}}
\newcommand{\examdate}{\placeholder{Exam date}}
\newcommand{\examduration}{\placeholder{duration of the exam}}
\newcommand{\problemfile}{dma-problems-april-2026} % problem file in this folder, without .tex
% ========================================

\usepackage{amssymb}
\usepackage{amsmath}
\usepackage{verbatim}
\usepackage{graphicx}
\usepackage{enumitem}
\usepackage{xspace}
\usepackage{tikz}
\usetikzlibrary{arrows,arrows.meta,automata,backgrounds,calc,
  decorations.pathmorphing,decorations.pathreplacing,fit,matrix,patterns,
  positioning,shapes,shapes.geometric,trees}
\usepackage{calc}
\usepackage{fancybox}
\usepackage{clrscode3e}
\usepackage{totcount}
\usepackage{geometry}
\geometry{a4paper, margin=1in}

\pagenumbering{gobble}

% Number of problems: every \question is a \section.
\regtotcounter{section}

\newcounter{answerpage}
\setcounter{answerpage}{1}
\regtotcounter{answerpage}

\newcommand{\answerpage}{\addtocounter{answerpage}{1}%
\placeholder{Where to write the answer, e.g.\ page \theanswerpage\ of the answer document.}
}

\newcommand{\subquestion}[2]{\vspace{.3cm}\noindent\begin{Sbox}
    \begin{minipage}{16cm}\textbf{#1} \xspace {#2}\xspace\answerpage\end{minipage}\end{Sbox}\fbox{\TheSbox}\\
    \vspace{.1cm}
    \noindent
    }

\tikzset{
  treenode/.style = {align=center, inner sep=0pt, text centered,
    font=\sffamily,text=white},
  black/.style = {treenode, circle, white, font=\sffamily\bfseries, draw=black,
    fill=black, text width=1.5em},
  red/.style = {treenode, circle, white, draw=red,
    fill=red, text width=1.5em, very thick},
  white/.style = {treenode, rectangle, draw=none,
    minimum width=0.5em, minimum height=0.5em},
  gray/.style = {treenode, rectangle, draw=black,
    minimum width=0.5em, minimum height=0.5em}
}

\newcommand{\question}[1]{\section{#1\mcpercent}}
\newcommand{\multiplechoicestart}{}
\newcommand{\multiplechoiceend}{}
\newcommand{\optionsstart}{\begin{enumerate}}
\newcommand{\optionsend}{\end{enumerate}}
\newcommand{\answer}{\item}
\newcommand{\itxanswer}[1]{} % For ITX, ignore in LaTeX
\newcommand{\qtipoints}[1]{} % For QTI/WISEflow, ignore in LaTeX
\newcommand{\oneanswer}{There is exactly one correct answer, but you may select more than one answer if you are in doubt.\xspace}
\newcommand{\multipleanswers}{Select one or more correct answers.\xspace}
\newcommand{\mcpercent}{\xspace (4\%)}

\ifanswers
\newcommand{\correct}{\checkmark}
\newcommand{\withanswers}{{\bf Problems with suggested answers}\\}
\newcommand{\explanation}[1]{{\bf Answer: }{#1}}
\newcommand{\pageshift}{}
\else
\newcommand{\correct}{}
\newcommand{\withanswers}{}
\newcommand{\explanation}[1]{}
\newcommand{\pageshift}{\newpage}
\fi

\begin{document}

\title{\Large \withanswers Exam in Discrete Mathematics and Algorithms\\
{\small \placeholder{Course code}}}

\author{\placeholder{Institution}}
\date{\examdate}

\maketitle

\section*{Instructions}

You have \examduration\ to solve the \total{section} problems described below.
\placeholder{How the grade is determined, e.g.\ from the total number of points.}
Each problem gives up to 4 points, so a perfect answer gives \the\numexpr 4*\totvalue{section}\relax\ points.

\medskip

\placeholder{Permitted aids, e.g.\ written material and locally installed software, and whether internet access and communication with others are allowed.}

\medskip

\placeholder{How the answers are submitted, e.g.\ the multiple-choice module of the exam system.}
Each multiple-choice question states either that there is exactly one correct answer or that there is a \emph{set} of one or more correct answers.
In both cases you may select more than one answer.
When multiple-choice problems are graded, wrong answers give negative points.
Examples:
\begin{itemize}
\item Problem X has 8 options, 4 of which are correct. Each check mark at a correct option gives 1 point, and each check mark at a wrong option gives $-1$ point. If, for example, you check all the correct options and one wrong option, you get 3 points.
\item Problem Y has 3 options, 1 of which is correct. A check mark at the correct option gives 4 points, and a check mark at a wrong option gives $-2$ points.
\end{itemize}
If you do not check any options you always get 0 points, and if you check all options you also always get 0 points.

\newpage

\input{\problemfile}

\end{document}
